\documentclass[11pt,a4paper]{article}
\usepackage[T1]{fontenc}
\usepackage[utf8]{inputenc}
\usepackage[french]{babel}
\usepackage{lmodern}
\usepackage{amsmath,amssymb}
\usepackage[margin=2.4cm]{geometry}
\usepackage{booktabs}
\usepackage{enumitem}
\usepackage{xcolor}
\usepackage{hyperref}
\hypersetup{colorlinks=true,linkcolor=blue!50!black}

\newcommand{\NW}{\mathrm{NW}}
\newcommand{\E}{\mathbb{E}}
\newcommand{\code}[1]{\texttt{#1}}

\title{Synthèse finale du programme M3\\
\large Livrable 06 (workspace \code{m3\_credit\_soc}) — réponses au protocole}
\author{Agent de recherche M3}
\date{5 juillet 2026}

\begin{document}
\maketitle

\begin{abstract}
M3 devait répondre à une question laissée ouverte par la double réplication
de M2 : un crédit \emph{productif} (conversion exclusive de liquidité en
capital), \emph{risqué} (pertes de stock et de flux) et \emph{topologique}
(réseau de contrats nominaux servis en liquidité) fait-il passer le système
d'un modèle démographique de renouvellement à un système
accumulation--relaxation ? Réponse d'ensemble : \textbf{non pour les
distributions, non pour la SOC — mais le crédit acquiert un rôle causal
démographique net (il divise par deux l'espérance de vie et la population)
et, découverte inattendue, il peut remplacer entièrement le plancher
d'absorption exogène $d_0$} : sans $d_0$, la dette contractuelle devient
l'unique mécanisme de mortalité et le système se borne quand même. Le
noyau Reed reste le seul générateur des formes ; l'hypothèse
Boltzmann--Gibbs échoue sur la liquidité elle-même ; le revenu est double
Pareto-lognormal partout (45 cellules de grille sur 45, 20 runs A/B sur
20). Chaque question du protocole reçoit ci-dessous une réponse explicite,
avec son degré de certitude.
\end{abstract}

\section{Réponses aux questions du protocole}

\subsection{Qu'est-ce qui est conservé / changé de M2 ?}
Conservé et vérifié : naissances Poisson, absence de cohorte fondatrice,
$d_0$, choc multiplicatif sans drift, extraction concave, dépréciation
réelle asymétrique du seul réel, taux géométrique (constitutif), marché
local par échantillons de taille $k$, cascade au point fixe,
reproductibilité par seed, fusion des contrats par paire (désormais par
défaut). Changé : état $(L, K)$ au lieu de $w$ ; partage de production
$s$ ; consommation $cL$ ; intérêts exigibles en liquidité (service
simultané au prorata) ; prêt $L_\ell \to K_b$ ; résidu de faillite en
nature ; avalanches causales tracées ; chocs corrélables. L'ancrage est
garanti par la propriété de réduction R1 (M3 dégénéré $=$ M2 bit à bit),
testée automatiquement : aucun écart A/B n'est un artefact
d'implémentation.

\subsection{Le mécanisme Reed suffit-il encore à tout expliquer ?}
\textbf{Oui, pour toutes les formes.} Corps de $\NW$ Fisk (9/10 baselines,
36/37 cellules de grille), queue effective $\alpha \approx 2{,}7$ stable
sur $20\,000$ pas, revenu double Pareto-lognormal \emph{partout} (10/10
baselines, 37/37 validations de grille — la signature Reed par excellence :
GBM $\times$ âges exponentiels), $K$ lognormal, réponse de $\alpha$ à
$\sigma$ identique à M2 ($3{,}7 \to 2{,}2$ pour $\sigma$ de $0{,}15$ à
$0{,}35$), insensibilité à $k$, $s$, $c$, $d_0$. L'ablation I le confirme
par l'absurde : sans renouvellement, extinction totale en $\sim 15$
espérances de vie — le mécanisme démographique est la condition
d'existence, pas un ingrédient parmi d'autres.

\subsection{La liquidité est-elle Boltzmann--Gibbs ?}
\textbf{Non — et c'est le verdict le plus instructif du programme.} $L$
est la variable pour laquelle l'argument de conservation locale était
formulé ; son corps est Fisk (l'exponentielle perd $\sim 700$ points
d'AIC ; $\mathrm{med/mean} = 0{,}93$ contre $\ln 2 = 0{,}69$ attendu),
queue légère. Le drift du corps de $L$ est pourtant quasi nul — la
dynamique est « équilibrée » — mais la forme n'est pas exponentielle,
parce que l'injection $(1-s)\alpha\sqrt{K_i}$ est hétérogène entre entités
et la fuite $cL$ multiplicative : les échanges conservatifs (intérêts,
$1{,}5\,\%$ de la production) sont un bruit de fond. Conclusion : dans
toute classe de modèles où la monnaie est un flux de production recyclé,
les hypothèses de Yakovenko--Rosser sont violées par construction ; le
cadre n'y est pas réfuté, il est \emph{sans domaine d'application}. La
dépendance à la conservation locale est donc confirmée en creux : pour
obtenir un corps exponentiel de $L$, il faudrait un régime dominé par les
transferts (échanges, dettes, défauts) et non par l'injection productive.

\subsection{$\NW$ a-t-il une structure à deux classes
(corps thermique $+$ queue superthermique) ?}
Non au sens de Yakovenko--Rosser. Le corps n'est ni exponentiel ni
« thermique » : son drift est ascendant ($+1{,}8$/pas, fenêtre fixe,
incréments à 1 pas), signature d'un régime de transport (réinjection en
bas, absorption aux bords), identique au diagnostic M2. La queue est
compatible Pareto à l'œil mais jamais identifiée face à la lognormale
tronquée (LR corrigé négatif dans toutes les fenêtres, tous les runs) ;
son exposant est stable mais il l'est \emph{aussi sans crédit}. La
description honnête de $\NW$ reste : une seule population Fisk-like
engendrée par le mélange démographique, sans frontière de classes.

\subsection{$K$ est-il stratifiant ou seulement mécanistique ?}
Mécanistique. Corps lognormal autour du point fixe
$K^\dagger = (s(1-\delta)\alpha/\delta)^2$, dans 10/10 baselines et sur
toute la grille ; le crédit ne le déforme pas (B identique), il ne change
même pas $K$ par tête (161 vs 155). $K$ ne stratifie pas la population :
c'est $\NW$ (via les créances/dettes et $L$) qui porte l'hétérogénéité, et
elle reste unimodale.

\subsection{Le crédit est-il causal ? (critère pré-enregistré)}
\textbf{Formellement oui ; substantiellement, il est causal sur la
démographie et neutre sur les distributions.} Le critère (deux signatures
modifiées par B) est atteint : population stationnaire ($\times 0{,}5$,
$> 20\sigma$), espérance de vie ($\times 0{,}5$), renouvellement du top
($0{,}985$ vs $0{,}864$, $> 5\sigma$), existence d'avalanches
multi-entités. Mais ces signatures sont les facettes d'un seul mécanisme —
mesuré par les ablations : \textbf{le canal létal est nominal} (C, crédit
non productif, reproduit tout l'effet : population $1055$ vs $1050$), les
pertes de créances n'y contribuent pas (D $=$ A), la perte de flux
d'intérêts pèse davantage (E récupère 40\,\% de l'écart), et l'effet
productif est nul en agrégat par tête (production/tête $11{,}2$ vs
$10{,}9$). Le crédit de M3 convertit de la liquidité sûre en capital
risqué assorti d'obligations : il \emph{accélère l'horloge de Reed} (moitié
moins de temps au-dessus du plancher) sans en changer la loi. Les
distributions — revenu, $L$, $\NW$, $K$, Gini, parts du top, exposants —
sont indiscernables de B (CCDF superposées, égalité des $\alpha$ au
centième).

\subsection{Le réseau compte-t-il au-delà du volume agrégé ?}
\textbf{Aucune évidence que oui.} L'ablation pré-enregistrée F a échoué par
conception (l'appariement aléatoire éteint le marché : volume $\div 6{,}5$
— échec documenté au livrable 05). La variante exploratoire F$''$
(prêteuse aléatoire, emprunteuse assortative) donne une population
intermédiaire à volume intermédiaire : à travers les trois règles
d'appariement, la population suit le \emph{volume} de façon monotone
(volume $76/41/12 \to$ population $1150/1665/2120$). Le pouvoir prédictif
individuel du réseau existe mais reste modeste (AUC levier $\to$ mort
$0{,}59$--$0{,}64$) et la concentration des créances est \emph{négativement}
corrélée aux faillites futures. Rien, dans ces données, ne nécessite la
topologie comme variable explicative au-delà de l'exposition agrégée.

\subsection{Les cascades sont-elles causales ou coïncidentes ?}
Le programme a corrigé la mesure de M2 (avalanches $=$ composantes du
graphe de pertes, racines classées par cause) : les cascades causales
\emph{existent} (perte de créance $\to$ insolvabilité du créancier),
mais elles sont petites (max 7 en baseline, $2\,\%$ des événements) et
\emph{sous}-dispersées (var/moy $= 0{,}06$). Sous chocs sectoriels, elles
grossissent (max 72 à $\rho_s = 0{,}8$) par synchronisation des défauts —
c'est de la causalité réelle, mais \emph{forcée de l'extérieur} : la
sur-dispersion suit la corrélation imposée continûment
($0{,}07 \to 0{,}38 \to 0{,}84$) sans jamais franchir le seuil de Poisson.

\subsection{Y a-t-il une signature SOC robuste ?}
\textbf{Non, sur aucun des trois critères pré-enregistrés.} (i) Pas de
dynamique accumulation--relaxation : la corrélation
concentration $\to$ faillites futures est systématiquement négative
($-0{,}2$ à $-0{,}6$) — la concentration suit les périodes calmes.
(ii) Pas d'avalanches à queue non triviale : var/moy $< 1$ partout, y
compris sous forçage corrélé maximal. (iii) Pas de zone critique sur la
grille : toutes les réponses ($\alpha$, population, avalanches) sont
monotones et lisses en $\sigma$, $k$, $s$, $c$, $d_0$, $\rho_s$ — aucune
transition, aucun plateau. Le système amplifie proportionnellement le
forçage qu'on lui impose ; il ne s'organise pas vers la criticité. Après
M2 (SOC « sans objet ») et M3 (SOC « réfutée dans le régime accessible »),
la charge de la preuve est écrasante : cette classe de modèles
(dette perpétuelle au taux marginal, liquidité produite en continu)
n'engendre pas d'auto-criticité.

\subsection{Le canal de fragilité nominal/réel fonctionne-t-il ?}
Il est implémenté et testé, mais \textbf{il ne s'active pratiquement
jamais} : 0 défaut de liquidité en baseline et dans presque toute la
grille ; 4 événements à $s = 0{,}9$ ; seule l'ablation H ($d_0 = 0$, levier
relatif élevé) en produit en nombre (71). Cause structurelle : au taux
marginal concurrentiel, le service de la dette coûte $\le 1{,}6\,\%$ de la
production ; la liquidité, alimentée par la production courante, couvre
toujours. La fragilité de type Minsky exige un ingrédient que M3 s'est
interdit : principal exigible, taux hors-équilibre, ou liquidité rare.

\subsection{Découverte positive : le plancher d'absorption peut être
endogène (ablation H)}
Sans $d_0$, seules les endettées peuvent mourir ($D >$ actifs) — et le
système se borne quand même, quasi-stationnairement, à une population
$\approx 5\times$ la baseline ($5574/5566/5649$ sur 3 seeds,
morts/naissances $= 0{,}98$ dans les trois cas, 67--79 défauts de
liquidité par run : le verdict est répliqué). Le crédit y est
\emph{constitutif} de la démographie : c'est le seul régime du programme où
retirer le crédit changerait qualitativement l'existence même du régime.
Ce résultat n'était pas anticipé par le protocole et désigne le M4 le plus
intéressant (cf. §3).

\subsection{Ce qui, de M2, est validé / à abandonner}
Validé et renforcé : le triplet générateur Reed, la prudence Pareto (le LR
corrigé ne s'inverse jamais), l'outillage statistique corrigé, la fusion
par paire, le taux géométrique constitutif, l'extensivité en $\lambda$
($N^*/\lambda = 114$--$115$ pour $\lambda \in \{3, 10, 30\}$). À
abandonner : l'espoir qu'un crédit comptablement actif suffise à créer de
la structure distributionnelle ; l'identification de \emph{toute} variable
de ces modèles à la variable Boltzmannienne ; l'ablation F par appariement
aléatoire (design invalide) ; le maintien de $\sigma^2$ total constant
comme seule convention pour les chocs corrélés (confond corrélation et
raréfaction du crédit — G1 inexploitable).

\section{Verdict d'ensemble}

M3 est un résultat négatif propre sur sa question principale, avec deux
acquis positifs nets : (1) le crédit a désormais un effet causal mesuré,
quantifié et mécanistiquement décomposé (nominal $\gg$ flux $>$ stock ;
productif $\approx 0$) — il raccourcit les vies sans changer les lois ;
(2) le plancher d'absorption, ingrédient le plus artificiel du cadre Reed,
peut être rendu endogène par la dette. La thèse Boltzmann--Pareto--SOC,
elle, sort du programme affaiblie sur ses trois volets : pas de corps
exponentiel (sur aucune variable, y compris la bonne), pas de Pareto
identifiée (partout dominée par la lognormale tronquée), pas de SOC (ni
accumulation-relaxation, ni avalanches critiques, ni zone robuste).

\section{M4 minimal recommandé}

Un seul changement structurel, dicté par H : \textbf{remplacer $d_0$ par
une dette de subsistance contractée à la naissance auprès d'entités
réelles} (pas d'un puits exogène), servie en liquidité. Le plancher
devient un réseau de créances : chaque mort inflige une perte réelle à des
créancières identifiées, le levier est élevé par construction (les
nouveau-nées naissent au plancher), et les défauts de liquidité existent
déjà dans ce régime (71 observés). Compléments optionnels si et seulement
si le canal reste faible : principal amortissable (l'échéancier crée
l'illiquidité), et le tampon $\beta_L$ comme variable de comportement.
Tout le reste de M3 (noyau Reed, outillage, protocole d'ablation corrigé
— F remplacée par une randomisation à volume contraint, G par corrélation
à dispersion idiosyncratique constante) est conservé. Prédiction réfutable
à pré-enregistrer pour M4 : dans le régime à plancher endogène, l'ablation
du crédit doit changer non seulement les niveaux mais l'\emph{existence}
du régime stationnaire ; si en outre les avalanches y restent
sous-Poisson, la thèse SOC devra être abandonnée pour cette famille de
modèles.

\section{Traçabilité}

50 tests verts (dont R1 : réduction bit à bit vers M2) ; 124 runs archivés
avec configuration, séries, instantanés, réseau et journal d'avalanches ;
\code{validation.json} par run, sans pooling inter-seed ; budget de
simulation consommé $20{,}3$ h sur 48, réplication H comprise (les trois
seeds de H : $10{,}1$ h à elles seules) ; journal de recherche
\code{NOTES.md} (leçons et décisions datées). Rapports : 01 (synthèse M2), 02 (spécification pré-enregistrée),
03 (implémentation), 04 (validation), 05 (résultats négatifs — aucun
n'est dilué ailleurs), 06 (le présent document).

\end{document}
